\sscp{Directions for further comparative study}{15-07-01}
In addition to needing more corpus-based grammars
and dictionaries, next steps for comparative studies in the region
probably need to track splits
and mergers at a lexically specific level with specific attention
to the apicals, as in central Maluku (see \crf{ch:CM} and \crf{ch:AS}).
For example, all \tsc{Sula-Buru}
and \tsc{\ili{Ambon}-Seram} languages show *j > \it{l} in *pajay, *pija.
All \tsc{Sula-Buru} and \tsc{\ili{Ambon}-Seram} languages show *j > {\0} in *qaləjaw.
PMP *j > {\0} in \tsc{Sula-Buru} in *ŋajan ‘name’,
*huaji ‘younger sibling (same sex)’,
but is preserved as *j > **l in \tsc{\ili{Ambon}-Seram} languages in the same words.
Data are limited or patterns are unclear for subgrouping purposes
for: *suja ‘pitfall or trail spikes made of sharpened bamboo’,
*quləj ‘maggot, caterpillar, larva of a metamorphosing insect’,
*pusəj ‘navel, umbilicus; mid-point’,
*laləj ‘housefly’, *waRəj ‘vine,
creeper’, *maja ‘dry up, evaporate’,
*hajək ‘smell, sniff, kiss’, *bujəq ‘foam, bubbles,
lather, scum, froth’, *qajəŋ ‘charcoal’.